% ID: SEL-D0026-OA-E02
% ESTADO: cotejada. Segunda pasada visual completa y borrador compilado revisados.
% Fuente: Selectividad/matematicas_aplicadas/suplemente_septiembre_EXAMEN 4-ANDALUCÍA.pdf
% SHA-256: c41ea8ffe46f2de3a4d7f82e43def97322e29653182ad3df54a57e2ead2beb19
% Página: 1; opción A; ejercicio original 2.

Sea la función
\[
f(x)=\begin{cases}
\dfrac{1}{a}x^2+1 & \text{si }x\leq 2\\[4pt]
-x+a & \text{si }x>2
\end{cases},\quad\text{con }a>0.
\]
\begin{enumerate}[label=\alph*)]
\item \textbf{(1.3 puntos)} Calcule el valor del parámetro $a$ para que la función sea continua en su dominio. En este caso, \textquestiondown{}sería derivable en su dominio?
\item \textbf{(1.2 puntos)} Para el valor $a=4$, represente gráficamente la función y halle la ecuación de la recta tangente a la gráfica de la función en el punto de abscisa $x=-1$.
\end{enumerate}
